% SUDS Showcase poster, August 14, 2026
% 4 ft x 3 ft landscape (121.92 cm x 91.44 cm), beamerposter.
% Build:  cd poster_v2 && pdflatex suds_poster_v4.tex
% Figures: cd poster_v2 && ../empirical/.venv/bin/python build_figures.py
%
\newcommand{\AuthorOne}{Yasir Alsugair}
\newcommand{\AuthorTwo}{Chuxuan Ai}
\newcommand{\ProjectName}{Hamiltonian Optimization and Sampling for Deep Learning}

\documentclass[final]{beamer}
\usepackage[size=custom,width=121.92,height=91.44,scale=1.32]{beamerposter}
\usepackage[T1]{fontenc}
\usepackage{lmodern}
% big operators come from the extension font at its 10pt design size, so without
% this \int and \sum stay small while everything else scales up to poster size
\usepackage{exscale}
\renewcommand{\familydefault}{\sfdefault}
% beamer's default sans math has no sans Greek, so it silently mixes a serif
% theta into sans variables. One serif math family throughout instead.
\usefonttheme[onlymath]{serif}
\usepackage{graphicx}
\usepackage{amsmath}
\usepackage{booktabs}
\usepackage{qrcode}
\graphicspath{{artifacts/figures/}{figures/}}

\setbeamertemplate{navigation symbols}{}
\setbeamersize{text margin left=2.5cm, text margin right=2.5cm}

\definecolor{uoftblue}{HTML}{002A5C}
\definecolor{inktext}{HTML}{20222B}
\definecolor{panelbg}{HTML}{F4F5F8}
\definecolor{softgray}{HTML}{8A8D99}
\definecolor{rtgold}{HTML}{D99A1B}
\definecolor{rtblue}{HTML}{1F6FB4}
\definecolor{hmcred}{HTML}{D1495B}

% gold accent rule marking the load-bearing panels
\newcommand{\dd}{\mathrm{d}}

\newcommand{\keyrule}{\textcolor{rtgold}{\rule{\linewidth}{2mm}}\\[0.5ex]}

% boxed key equations: \boxed inherits \fboxsep, whose 3pt default crowds a tall
% fraction against the rule at poster scale
\newcommand{\eqbox}[1]{{\setlength{\fboxsep}{9pt}\boxed{#1}}}

% Header band shares its outer edges with the block bars below, which run
% \bandmargin .. \paperwidth-\bandmargin on this geometry. Both the band and
% \linewidth are centred on the paper, so centring a \bandwidth box inside
% \linewidth lands on those edges without depending on beamer's column widths.
\newlength{\bandmargin}\setlength{\bandmargin}{2.13cm}
\newlength{\bandwidth}\setlength{\bandwidth}{\paperwidth}
\addtolength{\bandwidth}{-2\bandmargin}
\newcommand{\bandrule}{%
  \noindent\hbox to \linewidth{\hss
    \textcolor{uoftblue}{\rule{\bandwidth}{1.8mm}}%
  \hss}}

% Three logos on one line. The band gets an explicit height so what follows it
% does not depend on how deep the centred title box happens to be, and the
% flanking logo groups are set in zero-width boxes so the title stays centred on
% the paper whatever those groups measure.
\newlength{\bandheight}\setlength{\bandheight}{7.6cm}
\newcommand{\headerband}[1]{%
  \noindent\hbox to \linewidth{\hss
    \begin{minipage}[c][\bandheight][c]{\bandwidth}\leavevmode#1\end{minipage}%
  \hss}}

\setbeamercolor{normal text}{fg=inktext}
\setbeamercolor{block title}{fg=white,bg=uoftblue}
\setbeamercolor{block body}{fg=inktext,bg=panelbg}
\setbeamerfont{block title}{size=\Large,series=\bfseries}
\setbeamercolor{itemize item}{fg=uoftblue}

\setbeamertemplate{block begin}{%
  \begin{beamercolorbox}[ht=3.6ex,dp=1.2ex,leftskip=1.2ex]{block title}%
    \usebeamerfont*{block title}\insertblocktitle
  \end{beamercolorbox}%
  \nointerlineskip
  \begin{beamercolorbox}[leftskip=1.2ex,rightskip=1.2ex,vmode]{block body}%
    \vspace{0.9ex}\raggedright
    \setlength{\parskip}{0.21\baselineskip}\setlength{\parindent}{0pt}
}
\setbeamertemplate{block end}{%
  \vspace{0.9ex}\end{beamercolorbox}\vspace{0.24cm}
}

% centre a logo on the current baseline, so logos of unequal height sitting side
% by side align on their middles rather than their bottoms
\newcommand{\vlogo}[1]{\raisebox{-0.5\height}{#1}}

% a logo, or a labeled placeholder box if the file is missing
\newcommand{\logoslot}[4]{% file, label, height, placeholder width
  \IfFileExists{figures/#1}%
    {\includegraphics[height=#3]{#1}}%
    {\fcolorbox{softgray}{white}{\parbox[c][#3][c]{#4}{\centering
       \textcolor{softgray}{\small #2\\ {\tiny drop \texttt{\detokenize{figures/#1}}}}}}}}

% same idea for figures: a labeled placeholder rather than a failed build
\newcommand{\figslot}[3]{% file, width fraction, label
  \IfFileExists{artifacts/figures/#1}%
    {\includegraphics[width=#2\linewidth]{#1}}%
    {\IfFileExists{figures/#1}%
      {\includegraphics[width=#2\linewidth]{#1}}%
      {\fcolorbox{softgray}{white}{\parbox[c][6cm][c]{#2\linewidth}{\centering
         \textcolor{softgray}{\small #3\\ {\tiny missing \texttt{\detokenize{#1}}}}}}}}}

\begin{document}
\begin{frame}[t]

% ---------------- header band ----------------
% top margin matched to the 2.13 cm side margin
\vspace{2.0cm}
\headerband{%
  \rlap{\vlogo{\logoslot{dept_logo_trim.png}{U of T Statistical Sciences signature}{3.75cm}{12cm}}}%
  \parbox[c]{\bandwidth}{\centering
    {\textcolor{uoftblue}{\textbf{\LARGE Toward Scalable Bayesian Uncertainty for Machine Learning\\[0.2ex] with the Ray Tracing Sampler}}}\\[0.25cm]
    {\Large \AuthorOne{} \qquad \AuthorTwo}\\[0.18cm]
    {\large Supervisors: Prof.\ Joshua Speagle and Prof.\ Ricardo Baptista}}%
  \llap{%
    \vlogo{\logoslot{dsi_logo_standalone.png}{Data Sciences Institute logo}{2.5cm}{12cm}}%
    \hspace{3.0cm}%
    \vlogo{\logoslot{kaust_academy_logo_trim.png}{KAUST Academy logo}{6.0cm}{12cm}}}%
}
\vspace{0.22cm}
\bandrule
\vspace{0.18cm}

% ---------------- three columns ----------------
\begin{columns}[t]

% ============ column 1: why sample, where it lives, how you draw ============
\begin{column}{0.31\paperwidth}

\begin{block}{Why sample a neural network?}
\begin{center}
\figslot{fig0_data_bands.pdf}{0.60}{posterior draws vs one Adam fit}
\end{center}
\vspace{0.1cm}
\textbf{Ordinary training:} one set of weights $\rightarrow$ one function.\\
\textbf{Bayesian sampling:} a distribution over weights $\rightarrow$ uncertainty
over functions.
\vspace{-0.15cm}
\[
  \pi(w) \;=\; p(w \mid \mathcal{D}) \;\propto\;
  \underbrace{p(\mathcal{D} \mid w)}_{\text{fit}}\;
  \underbrace{p(w)}_{\text{prior}}
\]
\vspace{-0.45cm}

Every gold curve fits the data we have. They only disagree out in the shaded
gaps, and that disagreement is the uncertainty we are after.
\end{block}

\begin{block}{Where the probability lives}
\vspace{-0.15cm}
\[
  \mathbb{E}[f] \;=\; \int f(\theta)\, p(\theta)\, \mathrm{d}\theta
  \;=\; \int f(r)\, \textcolor{rtblue}{p(r)}\,
        \textcolor{rtgold}{r^{D-1}}\, \mathrm{d}r
\]
\begin{center}
\figslot{typical_set.pdf}{0.98}{density $\times$ volume, and the shell at $D=100$}
\end{center}
\vspace{0.1cm}
\textcolor{rtblue}{\textbf{Density}} peaks at the mode.
\textcolor{rtgold}{\textbf{Volume}} grows as $r^{D-1}$, for $D$ parameters.\\
Their product is the \emph{typical set}: a thin shell at
$\|\theta\|^2 \approx D$, relative width $\sqrt{2/D}$.

A sampler has to stay on that shell. If something quietly adds energy the chain
drifts outward, and none of the usual convergence diagnostics complain.
\end{block}

\begin{block}{MCMC in the era of big data}
You cannot draw from $\pi$ directly. A chain can end up there:
\vspace{-0.2cm}
\[
  w^{(0)} \rightarrow w^{(1)} \rightarrow \cdots \rightarrow w^{(M)},
  \qquad \{w^{(m)}\}_{m \gg 0} \sim \pi .
\]
\vspace{-0.5cm}

Gradient samplers steer by $\nabla \ln \pi$; Hamiltonian Monte Carlo (HMC) is
the standard one. Computed exactly, every step reads the whole dataset:
\vspace{-0.2cm}
\[
  \nabla \ln \pi(w) = \sum_{i=1}^{N} \nabla \ln p(y_i \mid w) + \nabla \ln p(w).
\]
\vspace{-0.5cm}

A minibatch $B \ll N$ is the standard fix:
\vspace{-0.2cm}
\[
  \begin{aligned}
  \hat{g}_B &= \tfrac{N}{B}\textstyle\sum_{i \in B}
               \nabla \ln p(y_i \mid w) + \nabla \ln p(w),\\[2pt]
  \mathbb{E}[\hat{g}_B] &= \nabla \ln \pi(w),
    \qquad \operatorname{Cov}[\hat{g}_B] \propto 1/B .
  \end{aligned}
\]
\vspace{-0.45cm}

The mini-batch gradient is unbiased, but no individual step ever uses the right
one. How much that per-step error costs you is what we set out to measure.
\end{block}

\end{column}

% ===================== middle column =====================
\begin{column}{0.31\paperwidth}

\begin{block}{Why ray tracing tolerates noisy gradients}

\textbf{1.\ Noise speeds HMC up. It only rotates a ray.}

\textbf{HMC} adds the gradient onto its velocity. A random push lengthens that
velocity whichever way it points, and nothing takes the energy back: the chain
heats up, drifts \emph{outward} off the shell, and its error bars come out too wide.

\textbf{Ray tracing} treats the posterior as an optical medium. Setting the
refractive index to \(n(\theta)=\pi(\theta)^{1/(D-1)}\) makes rays bend toward
higher probability. The sampler carries a \emph{unit} direction of travel \(u\),
and the gradient error enters as \(\sigma\zeta\):
\vspace{-0.03cm}
\[
    \eqbox{\;
    \frac{\dd u}{\dd s}
    = P_u\,\bigg(\nabla \ln n(\theta) + \frac{\sigma}{D-1}\,\zeta \bigg),
    \qquad
    P_u = I - uu^\top . \;}
\]
\vspace{-0.45cm}

Both act only in the tangent plane, so they can only rotate \(u\) on the
constant-speed sphere.

\vspace{0.08cm}
\begin{center}
\figslot{fig2_noise_geometry.pdf}{0.50}{HMC heating versus RT rotation}
\end{center}

\vspace{0.04cm}

\textbf{2.\ Many noisy turns become angular diffusion.}

Write \(b(h,\sigma)\) for the bias the noise leaves in equilibrium, and
\(\sigma_c(h)\sim h^{-p}\) for the noise tolerated at step \(h\). Averaged over the
noise, each kick contributes \(\sigma^2h^2/2(D-1)^2\) of diffusion \(\Delta_S\)
across the sphere of directions, and a trajectory of length \(S\) holds \(S/h\)
of them:
\vspace{-0.12cm}
\[
  \frac{S}{h}\cdot\frac{\sigma^2h^2}{2(D-1)^2}\,\Delta_S
  \;=\; O(\sigma^2h)\,\Delta_S ,
\]
\vspace{-0.45cm}

the order in \(h\) that leaves HMC at \(p=\tfrac12\).

\vspace{0.10cm}

\textbf{3.\ The leading diffusion preserves equilibrium.}

\vspace{-0.62cm}
\[
  \eqbox{\;
  \begin{aligned}
    \rho_0(\theta,u) &= \pi(\theta)\,\mu_S(u),
      \qquad \Delta_S\rho_0 = 0\\[3pt]
    &\Longrightarrow\;\; b(h,\sigma)=O(\sigma^2h^2)+o(\sigma^2)\\[3pt]
    &\Longrightarrow\;\; \sigma_c(h)=\Omega(h^{-1}),\quad p\ge1
  \end{aligned}
  \;}
\]
\vspace{-0.45cm}

The cancellation is \(\Delta_S\rho_0=0\): diffusion cannot spread directions that
are already uniform, and clean RT transport preserves \(\rho_0\) too.

\textbf{Individual paths wander at \(O(\sigma^2h)\). The distribution they settle
into does not move at that order.}

\vspace{0.10cm}

\textbf{4.\ Is the bound tight?}

Theory gives \(p\ge1\) for ray tracing against \(p=\tfrac12\) for HMC, so halving
the step buys HMC \(1.4\times\) the tolerable noise and ray tracing at least
\(2\times\).

\vspace{0.06cm}
\begin{center}
\figslot{fig8_dimension_scaling.pdf}{0.87}{tolerance and exponent against dimension}
\end{center}
\vspace{0.06cm}
{\small The bias is read against the clean chain at the \emph{same} step, so
integration error cancels:
\(b=\big[\mathbb{E}_{h,\sigma}\|\Theta\|^2-\mathbb{E}_{h,0}\|\Theta\|^2\big]/D
=(\sigma/\sigma_c)^{2}+o(\sigma^2)\).}



\end{block}



\end{column}

% ============ column 3: the measurement, then the real test ============
\begin{column}{0.31\paperwidth}

\begin{block}{Failure modes, measured}

Injected noise, \(\sigma = 20\), \(D = 100\): stochastic-gradient HMC
equilibrates three times off the shell with split \(\hat R = 1.00\), no alarm;
ray tracing holds the shell.
\begin{center}
\figslot{shell_measured_poster.pdf}{0.68}{the shell, measured: off-shell chains still pass $\hat R$}
\end{center}
\vspace{0.04cm}

The theory assumes the gradient error is independent of position and symmetric
in sign. Force both onto a real minibatch (\emph{idealized} below) and the
exponent holds; leave the batch as it comes and it decays with every halving, so
no single \(p\) survives.

\vspace{0.02cm}
\begin{center}
\figslot{fig9_erosion_and_frontier.pdf}{0.72}{exponent erosion beside the compute frontier}
\end{center}
\vspace{0.00cm}

\end{block}

\begin{block}{Uncertainty quality on the Gaia posterior}
Gaia XP regression, 126{,}156 stars, \(D = 11{,}394\), sampled with minibatch
ray tracing. Key: \textbf{\textcolor{softgray}{point (MAP)}},
\textbf{\textcolor{rtblue}{deep ensemble}},
\textbf{\textcolor{hmcred}{SGHMC}},
\textbf{\textcolor{rtgold}{ray tracing}}.
\begin{center}
\figslot{exp7_structure.pdf}{0.66}{per-group calibration, z std per label-error bin}\\[0.04cm]
\figslot{exp7_members.pdf}{0.66}{test NLL against members averaged}
\end{center}
\end{block}

\begin{block}{Conclusion}
\begin{itemize}\setlength{\itemsep}{0.02cm}
  \item Gradient noise can only rotate a unit ray, and the fitted step
    exponent confirms the bound: \(p=1\) at every dimension tested
    (\(D=16\) to \(1024\)), against \(p=\tfrac12\) for stochastic HMC.
  \item Real minibatches break the noise assumptions and erode the clean
    exponent, yet the compute advantage survives (\(3.3\times\) cheaper per
    effective sample; stochastic HMC \(5.1\times\) dearer). Still open:
    whether \(p=1\) is exact in stationarity.
  \item On a real survey posterior, only ray tracing keeps per-group
    calibration flat, and member averaging passes a deep ensemble on test
    NLL.
\end{itemize}
\end{block}

\begin{block}{Acknowledgments}
The \ProjectName{} project is supported by the Data Sciences Institute,
University of Toronto. We thank the KAUST Academy for its support.
\end{block}

\end{column}

\end{columns}

\end{frame}
\end{document}
